\section{Síntesis y ampliación}

\subsection{Repaso de los puntos clave}

DDIM es un trabajo fundacional para la aceleración de la muestreo en modelos de difusión, cuyos principales aportes se pueden resumir en los siguientes aspectos:

\begin{enumerate}
    \item \textbf{Proceso anterior no markoviano}: Al definir directamente la distribución de salto y la distribución a posteriori (en lugar de la distribución de transición anterior), se construyó un marco de difusión más general que DDPM. La distribución a posteriori introduce un parámetro libre $\sigma_t$ para controlar el nivel de aleatoriedad en el muestreo.

    \item \textbf{Invariancia del ELBO}: Dado que la distribución de salto se mantiene consistente con DDPM, el objetivo de entrenamiento de DDIM es completamente equivalente al de DDPM. Esto significa que los modelos existentes de DDPM pueden utilizarse directamente para la muestreo con DDIM sin necesidad de ningún reentrenamiento.

    \item \textbf{Muestreo determinista}: Cuando $\sigma_t = 0$, el proceso de muestreo se vuelve completamente determinista, estableciendo una correspondencia biyectiva entre el espacio de ruido y el espacio de datos. Esto aporta reproducibilidad, un espacio latente semántico y conexiones con solucionadores de EDO (ecuaciones diferenciales ordinarias).

    \item \textbf{Generación acelerada}: Mediante estrategias de muestreo por subsecuencias, se puede reducir el número de pasos desde 1000 hasta aproximadamente 50, logrando una aceleración de 20 veces. El modo de muestreo determinista muestra mayor robustez al reducirse el número de pasos.

    \item \textbf{Operaciones en el espacio latente}: La simetría codificación-decodificación de DDIM determinista hace posibles operaciones semánticas como la interpolación en el espacio latente.
\end{enumerate}

\begin{importantbox}{El impacto profundo de DDIM}
DDIM no es solo un método práctico de aceleración de muestreo por sí mismo, sino que es más importante revelar un principio de diseño clave en los modelos de difusión: \textbf{el desacoplamiento entre entrenamiento e inferencia}. Una vez entrenado el modelo, durante la inferencia se pueden elegir libremente diferentes estrategias de muestreo (número de pasos, nivel de aleatoriedad, solucionadores EDO, etc.). Este principio ha sido heredado por casi todo el trabajo posterior en aceleración de modelos de difusión, incluidos DPM-Solver, Consistency Models y Progressive Distillation.
\end{importantbox}

\subsection{Direcciones futuras}

Basado en el marco de DDIM, la investigación posterior avanzó desde múltiples direcciones para mejorar la eficiencia del muestreo en modelos de difusión:

\begin{enumerate}
    \item \textbf{Solucionadores EDO de orden superior} (DPM-Solver, DPM-Solver++): Al considerar el DDIM determinista como un método de Euler para una EDO de flujo probabilístico, el uso de solucionadores de orden superior (2.º, 3.er orden) permite obtener mayor precisión al mismo número de pasos o alcanzar la misma calidad con menos pasos.
    \item \textbf{Marco continuo en tiempo para SDE/ODE} (Score SDE): Al extender DDPM/DDIM de tiempo discreto a tiempo continuo, se unifican los modelos basados en puntajes y los modelos de difusión.
    \item \textbf{Distilación de conocimiento} (Progressive Distillation, Consistency Distillation): Mediante la distilación, se reduce aún más el número de pasos hasta 1-4.
    \item \textbf{Flow Matching}: Se rediseñan las trayectorias de difusión desde una perspectiva de transporte óptimo para lograr trayectorias de muestreo más rectas y menos pasos.
\end{enumerate}

\subsection{Lecturas adicionales}

\begin{itemize}
    \item Song, J., Meng, C., \& Ermon, S. (2021). Denoising Diffusion Implicit Models. \textit{ICLR 2021}. \href{https://arxiv.org/abs/2010.02502}{arXiv:2010.02502}
    \item Song, Y., Sohl-Dickstein, J., Kingma, D. P., Kumar, A., Ermon, S., \& Poole, B. (2021). Score-Based Generative Modeling through Stochastic Differential Equations. \textit{ICLR 2021}. \href{https://arxiv.org/abs/2011.13456}{arXiv:2011.13456}
    \item Ho, J., Jain, A., \& Abbeel, P. (2020). Denoising Diffusion Probabilistic Models. \textit{NeurIPS 2020}. \href{https://arxiv.org/abs/2006.11239}{arXiv:2006.11239}
    \item Lu, C., Zhou, Y., Bao, F., Chen, J., Li, C., \& Zhu, J. (2022). DPM-Solver: A Fast ODE Solver for Diffusion Probabilistic Model Sampling. \textit{NeurIPS 2022}. \href{https://arxiv.org/abs/2206.00927}{arXiv:2206.00927}
    \item Song, Y., \& Dhariwal, P. (2023). Consistency Models. \textit{ICML 2023}. \href{https://arxiv.org/abs/2303.01469}{arXiv:2303.01469}
\end{itemize}

%% --- Fin del contenido --- %%

\end{document}
